%% 中英文摘要

\cnabstract

无穷级数是数学分析的核心内容之一，其求解问题在理论研究与工程应用中具有重要意义。本文以无穷级数的解及其存在形式为研究对象，系统探讨了解析解与数值解两种求解路径，并提出了一种谱系化理解框架，旨在为无穷级数的求解提供统一的分类视角。

首先，本文回顾了无穷级数的基本概念与收敛判别理论，明确了解析解与数值解的定义及其界定标准。在解析解研究方面，本文梳理了等比级数求和、裂项求和、泰勒展开及傅里叶级数展开等主要求解方法。在数值解研究方面，本文讨论了直接部分和逼近法、欧拉变换加速、欧拉-麦克劳林积分修正、帕德逼近以及基于 Wynn $\varepsilon$-算法和 Richardson 的序列外推等数值加速算法，并对不同方法在截断误差与舍入误差之间的平衡关系进行了分析。

在此基础上，本文构建了无穷级数解的谱系化框架，指出无穷级数的解仅有解析解与数值解两种基本形式，并讨论二者的内在关联：数值解可以解析解为基础，利用解析结构构造更有效的数值算法。为验证该框架的合理性，本文结合数值实验，对等比级数、交错调和级数、$\zeta(3)$ 常数及幂级数等典型实例进行了计算分析。误差分析与收敛变化图表表明，以解析解为基础的加速算法能够有效提高收敛效率，同时也反映了不同方法在计算机有限字长限制下截断误差与舍入误差的变化规律。

本文的研究表明，谱系化视角及数值实验验证有助于加深对无穷级数解的存在形式的理解，并可为数学分析教学与科学计算应用提供一种系统化的参考思路。

\cnkeywords{无穷级数；解析解；数值解；收敛性；谱系化}

\enabstract

Infinite series is one of the core topics in mathematical analysis, and its solution problem holds significant importance in both theoretical research and engineering applications. This thesis systematically investigates analytical and numerical solutions of infinite series, proposing a spectralized understanding framework aimed at providing a unified classification perspective.

First, the thesis reviews the fundamental concepts and convergence criteria of infinite series. In the study of analytical solutions, major methods including geometric series, telescoping series, Taylor expansion, and Fourier series expansion are systematically examined. In the study of numerical solutions, the thesis focuses on modern high-efficiency acceleration algorithms such as direct partial sum approximation, Euler transform, Euler-Maclaurin integral correction, Padé approximants, and sequence extrapolations based on the Wynn $\varepsilon$-algorithm and Richardson method, providing an in-depth analysis of their limitations regarding truncation and underlying round-off errors.

On this basis, a spectralized framework is constructed, clarifying that infinite series solutions exist in only two fundamental forms---analytical solutions and numerical solutions---while revealing their intrinsic connection: numerical solutions can be built upon analytical solutions, leveraging analytical structures to construct more efficient numerical algorithms. To validate this framework, the thesis incorporates numerical experiments on typical instances including geometric series, alternating harmonic series, $\zeta(3)$ constant, and power series. The error analysis and convergence charts indicate that acceleration algorithms grounded in analytical structures can significantly improve convergence efficiency, while also revealing the evolution of truncation and round-off errors under the finite word-length limit of computers.

The research demonstrates that the spectralized perspective combined with numerical verification contributes to deepening the understanding of the existence forms of infinite series solutions, providing a systematic approach for mathematical analysis teaching and scientific computing applications.

\enkeywords{infinite series; analytical solution; numerical solution; convergence; spectralization}

\clearpage
